\documentclass{article}
\usepackage[margin=1in]{geometry}
\usepackage{longtable}
\usepackage{hyperref}
\usepackage{array}
\usepackage{booktabs}
\usepackage{enumitem}

\title{Software Requirements Specification (SRS)\\
\large LA City Sidewalk Assessment Simulation}
\author{James Barron --- SRS Lead / Requirements}
\date{\today}

\begin{document}

\maketitle
\newpage
\tableofcontents
\newpage


\section{Revision History}


\begin{longtable}{|p{1.3in}|p{1in}|p{2.6in}|p{0.8in}|}
\hline
\textbf{Name} & \textbf{Date} & \textbf{Reason for Change} & \textbf{Version} \\
\hline
James Barron & 2025-01-15 & Initial Snapshot 1 SRS draft & 1.0 \\
\hline
James Barron & 2025-04-10 & Snapshot 2: added data processing and displacement pipeline requirements & 2.0 \\
\hline
James Barron & 2025-07-08 & Snapshot 3: added semi-autonomous navigation requirements & 3.0 \\
\hline
James Barron & 2025-10-20 & Snapshot 4: final revision, completed and deferred features, future work & 4.0 \\
\hline
\end{longtable}


\section{Introduction}


\subsection{Purpose}
This Software Requirements Specification (SRS) defines the functional and nonfunctional requirements
for the LA City Sidewalk Assessment simulation system. The document is intended for use by the
development team, project advisor, QA testers, and LA City Bureau of Engineering (BOE) reviewers.
It serves as the single source of truth for what the system is required to do and the constraints
under which it must operate.

This SRS is a living document updated at each project snapshot. Version 1.0 establishes the project
scope and baseline requirements. Subsequent versions expand on technical features as they are
designed and validated.

\subsection{Product Scope}
The LA City Sidewalk Assessment system supports a rover-based sidewalk inspection workflow. The
software goal is to collect sidewalk scan data, process LiDAR and point-cloud information, identify
sidewalk joints and cracks, measure pavement displacement, and produce outputs suitable for LA City
Bureau of Engineering review. The simulated 12-month development timeline covers:

\begin{itemize}
  \item Field data collection via rover-mounted and iPhone LiDAR sensors.
  \item Point-cloud processing to generate orthoimages and Digital Elevation Models (DEMs).
  \item Deep-learning-based segmentation using a U-Net model to detect cracks and joints.
  \item Vertical and horizontal displacement measurement and CSV export.
  \item Semi-autonomous rover navigation with GPS waypoints, obstacle detection, and safety stops.
  \item A reviewer-facing dashboard for accessing assessment outputs.
\end{itemize}

\subsection{Definitions, Acronyms, and Abbreviations}
See Section~\ref{sec:glossary} for a full glossary.

\subsection{References}
\begin{itemize}
  \item CS3338 Final Project Assignment Description.
  \item LA City Bureau of Engineering sidewalk inspection guidelines.
  \item Americans with Disabilities Act (ADA) Accessibility Standards.
  \item CS3338 LA City Sidewalk Project Step-by-Step Plan.
\end{itemize}

\subsection{Overview}
Section~\ref{sec:overall} provides a high-level product description. Section~\ref{sec:interfaces}
describes external interfaces. Section~\ref{sec:functional} lists all functional requirements.
Section~\ref{sec:nonfunctional} covers nonfunctional requirements. Section~\ref{sec:legal} addresses
legal and ethical considerations. Section~\ref{sec:glossary} is the glossary.


\section{Overall Description}
\label{sec:overall}


\subsection{Product Perspective}
The LA City Sidewalk Assessment system is a new, self-contained simulation of a software product
that would be deployed by city engineering teams to replace or augment manual sidewalk inspection.
The system interfaces with physical sensors (LiDAR, GPS, cameras) through a Raspberry Pi / ROS
layer, and with cloud or local storage through a Python-based processing backend.

\subsection{Product Functions}
At a high level the system shall:
\begin{enumerate}
  \item Accept raw scan data from rover-mounted or iPhone LiDAR devices.
  \item Convert point-cloud files (LAS/PLY) into orthoimages and elevation images.
  \item Run a U-Net segmentation model to produce crack and joint masks.
  \item Compute vertical displacement from DEM and mask alignment.
  \item Research and document horizontal displacement methodology.
  \item Export results as CSV files and visualization images.
  \item Support semi-autonomous rover navigation via GPS waypoints.
  \item Detect obstacles and execute manual override or safe-stop procedures.
  \item Log telemetry data for later review.
  \item Present outputs through a browser-based reviewer dashboard.
\end{enumerate}

\subsection{User Classes and Characteristics}

\begin{longtable}{|p{1.5in}|p{2in}|p{2.5in}|}
\hline
\textbf{User Class} & \textbf{Description} & \textbf{Skill Level / Needs} \\
\hline
Rover Operator & Field technician who deploys the rover and monitors navigation. &
  Moderate technical skill. Needs a simple UI for mission start, stop, manual override, and live
  telemetry view. \\
\hline
Software Developer & Team member who builds, maintains, and extends the processing pipeline. &
  High technical skill. Needs clear API contracts, Docker documentation, and test coverage. \\
\hline
QA Tester & Team member responsible for test case creation and TestRail reporting. &
  Moderate technical skill. Needs traceability from requirements to test cases and clear pass/fail
  criteria. \\
\hline
BOE / Civil Engineering Reviewer & City engineer who reviews sidewalk assessment outputs. &
  Domain expert, low software skill. Needs readable CSV exports, visualizations, and a simple
  dashboard. \\
\hline
Project Advisor & CS professor who evaluates the simulation for academic correctness. &
  High academic skill. Needs complete documentation, Jira evidence, GitHub history, and TestRail
  reports. \\
\hline
\end{longtable}

\subsection{Operating Environment}
\begin{itemize}
  \item \textbf{Hardware:} Rover with Raspberry Pi 4, iPhone 14 Pro (LiDAR), GPS module, obstacle
    detection camera.
  \item \textbf{Software:} Ubuntu 22.04 on the processing workstation; ROS Noetic on the
    Raspberry Pi; Docker and docker-compose for service orchestration; Python 3.10 for the data
    processing and ML pipeline.
  \item \textbf{Network:} Local network or Wi-Fi for rover-to-workstation communication;
    internet access for GitHub, Jira, and TestRail.
  \item \textbf{Storage:} Local volume or Azure Blob Storage placeholder for raw scans,
    processed data, and model weights.
\end{itemize}

\subsection{Design and Implementation Constraints}
\begin{itemize}
  \item The system must run within Docker containers to ensure reproducibility.
  \item Machine learning inference must run on available CPU/GPU without cloud dependency.
  \item LiDAR data files must be in LAS or PLY format.
  \item The U-Net model weights file must be stored at \texttt{storage/models/unet\_model.h5}.
  \item All configuration must be done through environment variables; no hard-coded credentials.
\end{itemize}

\subsection{Assumptions and Dependencies}
\begin{itemize}
  \item It is assumed that field scan data is already collected before software processing begins.
  \item The team does not need to train the U-Net model from scratch; pre-trained weights are
    supplied.
  \item GPS accuracy is assumed sufficient for waypoint navigation within a city block.
  \item The Jira project and GitHub repository are set up by the Project Manager before any
    development work begins.
\end{itemize}


\section{External Interface Requirements}
\label{sec:interfaces}


\subsection{User Interfaces}
\begin{itemize}
  \item \textbf{Dashboard (UI-1):} A browser-based dashboard accessible at
    \texttt{http://localhost:5173}. Displays mission status, recent scan results, displacement
    summaries, and links to CSV exports.
  \item \textbf{Upload Scan (UI-2):} A file upload form accepting LAS and PLY point-cloud files
    up to 2 GB per scan.
  \item \textbf{Scan Review (UI-3):} A page showing orthoimage and segmentation mask
    side-by-side for a selected scan.
  \item \textbf{Displacement Results (UI-4):} A table and chart view of vertical displacement
    measurements per sidewalk segment.
  \item \textbf{Rover Navigation (UI-5):} A map-based panel showing live rover position,
    waypoint queue, obstacle events, and manual override controls.
  \item \textbf{Test Logs (UI-6):} A read-only log view for telemetry and processing events.
  \item \textbf{Admin Settings (UI-7):} A settings page for configuring storage paths, model
    path, and notification thresholds.
\end{itemize}

\subsection{Hardware Interfaces}
\begin{itemize}
  \item \textbf{Rover / Raspberry Pi (HW-1):} The backend communicates with the Raspberry Pi
    over a local network socket using ROS Bridge or a simulated REST interface.
  \item \textbf{iPhone LiDAR (HW-2):} The iPhone 14 Pro records point-cloud files that are
    transferred to the workstation via file export.
  \item \textbf{GPS Module (HW-3):} A NMEA-compatible GPS module attached to the rover provides
    latitude/longitude coordinates at 1 Hz.
  \item \textbf{Obstacle Camera (HW-4):} A forward-facing camera on the rover streams frames
    to the obstacle detection module.
\end{itemize}

\subsection{Software Interfaces}
\begin{itemize}
  \item \textbf{Python 3.10 (SW-1):} All processing scripts run on Python 3.10.
  \item \textbf{TensorFlow / Keras (SW-2):} Used to load and run the U-Net segmentation model.
  \item \textbf{ROS / ROS Bridge (SW-3):} Simulated or real ROS topics for rover control,
    telemetry, and navigation.
  \item \textbf{PostgreSQL 16 (SW-4):} Stores scan metadata, job status, and displacement
    records.
  \item \textbf{FastAPI (SW-5):} Backend REST API for file upload, job management, and result
    retrieval.
  \item \textbf{React / Vite (SW-6):} Frontend single-page application served on port 5173.
  \item \textbf{Docker / docker-compose (SW-7):} Orchestrates all services.
  \item \textbf{GitHub (SW-8):} Version control and pull-request workflow.
  \item \textbf{Jira (SW-9):} Task tracking and sprint management.
  \item \textbf{TestRail (SW-10):} Test case management and report export.
\end{itemize}

\subsection{Communication Interfaces}
\begin{itemize}
  \item \textbf{HTTP/REST (COM-1):} Frontend communicates with the backend over HTTP on port
    8000 using JSON payloads.
  \item \textbf{WebSocket (COM-2):} Live rover telemetry is streamed to the frontend over a
    WebSocket connection.
  \item \textbf{File System Volume (COM-3):} The backend and ML worker share a mounted Docker
    volume at \texttt{/app/storage} for large scan files and processed outputs.
\end{itemize}


\section{Functional Requirements}
\label{sec:functional}


\subsection{Snapshot 1 --- Baseline Requirements}

\begin{longtable}{|p{0.8in}|p{5.4in}|}
\hline
\textbf{ID} & \textbf{Requirement} \\
\hline
FR-1 & The system shall accept LAS and PLY point-cloud files uploaded through the web UI. \\
\hline
FR-2 & The system shall convert uploaded point-cloud data into a 2D orthoimage and a
  Digital Elevation Model (DEM) image. \\
\hline
FR-3 & The system shall run a pre-trained U-Net model on the orthoimage to produce a
  binary segmentation mask identifying sidewalk joints and cracks. \\
\hline
FR-4 & The system shall calculate vertical displacement measurements by aligning the
  segmentation mask with the DEM elevation values. \\
\hline
FR-5 & The system shall export displacement measurements and metadata as a CSV file. \\
\hline
FR-6 & The system shall store raw scan files, processed orthoimages, DEMs, masks, and
  CSV exports in a designated storage volume. \\
\hline
FR-7 & The system shall display processed results in the Displacement Results dashboard
  view. \\
\hline
FR-8 & The system shall maintain a processing job queue, showing status as pending,
  running, completed, or failed. \\
\hline
\end{longtable}

\subsection{Snapshot 2 --- Data Processing and Displacement Pipeline Requirements}

\begin{longtable}{|p{0.8in}|p{5.4in}|}
\hline
\textbf{ID} & \textbf{Requirement} \\
\hline
FR-9 & The system shall parse LAS 1.2--1.4 and PLY point-cloud files using the Laspy
  and Open3D libraries without loss of point attributes. \\
\hline
FR-10 & The system shall generate an orthoimage with a configurable ground sample distance
  (GSD), defaulting to 1 mm per pixel. \\
\hline
FR-11 & The system shall generate a DEM image aligned to the orthoimage resolution and
  coordinate frame. \\
\hline
FR-12 & The system shall validate that the orthoimage and DEM have matching dimensions and
  coordinate extents before running segmentation; a mismatch shall produce a logged error
  and halt the job. \\
\hline
FR-13 & The system shall apply noise removal and elevation smoothing (median filter,
  configurable kernel size) to the DEM before displacement computation. \\
\hline
FR-14 & The system shall compute vertical displacement as the signed elevation difference
  between adjacent sidewalk panel segments identified by the U-Net mask. \\
\hline
FR-15 & The system shall flag any sidewalk segment where vertical displacement exceeds
  a configurable threshold (default: 0.5 inches / 12.7 mm) as requiring remediation. \\
\hline
FR-16 & The system shall document horizontal displacement methodology in the SDD and
  SRS; full horizontal displacement calculation is deferred to a future release. \\
\hline
FR-17 & The system shall produce an annotated visualization image overlaying the
  segmentation mask and displacement values on the orthoimage. \\
\hline
FR-18 & The system shall record processing start time, end time, and error messages for
  each job in the PostgreSQL database. \\
\hline
FR-19 & The CSV export shall include columns: scan\_id, segment\_id, orthoimage\_path,
  dem\_path, mask\_path, vertical\_displacement\_mm, flag\_remediation,
  processed\_at. \\
\hline
FR-20 & The system shall validate model output confidence; if the U-Net produces an
  average mask confidence below 0.6 for a scan, the result shall be logged as
  low-confidence and highlighted in the dashboard. \\
\hline
\end{longtable}

\subsection{Snapshot 3 --- Semi-Autonomous Navigation Requirements}

\begin{longtable}{|p{0.8in}|p{5.4in}|}
\hline
\textbf{ID} & \textbf{Requirement} \\
\hline
FR-21 & The system shall accept a GPS waypoint list (latitude, longitude, optional label)
  uploaded via the Rover Navigation UI or through a JSON file. \\
\hline
FR-22 & The system shall execute waypoints sequentially; the rover shall proceed to each
  waypoint in the order specified and shall not skip waypoints unless a safety override
  is triggered. \\
\hline
FR-23 & The system shall stop the rover automatically and alert the operator when the
  rover's calculated path to the next waypoint is blocked by a detected obstacle. \\
\hline
FR-24 & The obstacle detection module shall process frames from the forward-facing camera
  at a minimum of 5 frames per second during navigation. \\
\hline
FR-25 & The system shall provide a manual override control in the UI that immediately
  suspends autonomous navigation and transfers control to operator joystick or directional
  input. \\
\hline
FR-26 & The system shall execute a safe-stop procedure (immediate motor halt, alarm log
  entry) if communication with the rover is lost for more than 3 consecutive seconds. \\
\hline
FR-27 & The system shall log telemetry at a minimum of 1 Hz during navigation, including:
  timestamp, latitude, longitude, heading, speed, battery level, and obstacle status. \\
\hline
FR-28 & The Rover Navigation UI shall display the rover's live GPS position on a map
  updated at least once per second. \\
\hline
FR-29 & The system shall display remaining battery level on the Rover Navigation UI and
  shall alert the operator when battery falls below 20\%. \\
\hline
FR-30 & The system shall store all telemetry logs in the database and allow export as
  CSV from the Test Logs UI. \\
\hline
\end{longtable}

\subsection{Snapshot 4 --- Final Requirements (Completed and Deferred)}

\begin{longtable}{|p{0.8in}|p{5.4in}|}
\hline
\textbf{ID} & \textbf{Requirement} \\
\hline
FR-31 & The system shall include a field testing report section in the reviewer dashboard
  summarizing results from at least one real or simulated sidewalk test run. \\
\hline
FR-32 & The system shall provide a user manual accessible from the dashboard that
  describes all features, file formats, Docker setup, and known limitations. \\
\hline
FR-33 & The system shall version all exported CSV and visualization files by scan ID
  and processing timestamp so older results are not overwritten. \\
\hline
FR-34 & [DEFERRED] Full horizontal displacement calculation using paired scan alignment
  is deferred to a future release pending additional LiDAR hardware calibration data. \\
\hline
FR-35 & [DEFERRED] Real-time on-rover inference (running U-Net on the Raspberry Pi)
  is deferred; the current design processes scans on the workstation after upload. \\
\hline
FR-36 & [DEFERRED] Cloud processing integration (Azure Blob / S3) is planned but not
  implemented in this simulation; local volume storage is used throughout. \\
\hline
\end{longtable}


\section{Nonfunctional Requirements}
\label{sec:nonfunctional}


\subsection{Performance}
\begin{itemize}
  \item The system shall process a single sidewalk scan (up to 100 MB point-cloud file)
    from upload through CSV export in under 5 minutes on the reference workstation
    (8-core CPU, 16 GB RAM, no GPU).
  \item The dashboard shall load within 3 seconds on a standard broadband connection.
  \item The Rover Navigation UI shall update the live GPS position within 2 seconds of
    receiving a telemetry packet.
\end{itemize}

\subsection{Safety}
\begin{itemize}
  \item The rover shall execute a safe-stop immediately upon loss of communication or
    manual override trigger (FR-26, FR-25).
  \item The system shall not command the rover to proceed through a flagged obstacle
    zone without explicit operator confirmation.
  \item Battery level alerts shall fire before the rover is stranded in the field
    (FR-29).
\end{itemize}

\subsection{Security}
\begin{itemize}
  \item No credentials shall be hard-coded in source files; all secrets shall be
    supplied via environment variables or Docker secrets.
  \item The admin settings page (UI-7) shall be accessible only to authenticated users.
  \item Field scan data containing images of pedestrians or private property shall be
    stored in access-controlled storage and not exposed publicly.
\end{itemize}

\subsection{Reliability and Maintainability}
\begin{itemize}
  \item All processing jobs shall be retryable; a failed job shall log the error and
    allow manual re-trigger without data loss.
  \item The system shall pass all TestRail regression cases before any snapshot merge
    into the main branch.
  \item Code shall follow consistent formatting and be documented with inline comments
    sufficient for a new developer to onboard within one day.
\end{itemize}

\subsection{Usability}
\begin{itemize}
  \item A BOE reviewer with no software background shall be able to access the
    dashboard, view displacement results, and export a CSV without consulting the user
    manual.
  \item All error messages displayed in the UI shall be written in plain English and
    shall suggest a corrective action.
\end{itemize}

\subsection{Portability}
\begin{itemize}
  \item The system shall run on any machine where Docker and docker-compose are
    installed, regardless of the host operating system.
  \item The processing scripts shall produce identical output on Windows (via WSL2),
    macOS, and Linux.
\end{itemize}


\section{Database and Storage Requirements}


\begin{itemize}
  \item The PostgreSQL database shall store: scans, jobs, displacement\_records,
    telemetry\_logs, and users tables.
  \item Raw scan files shall be stored on the Docker volume at
    \texttt{storage/raw/<scan\_id>/}.
  \item Processed outputs (orthoimage, DEM, mask, CSV, visualization) shall be stored
    at \texttt{storage/processed/<scan\_id>/}.
  \item Model weights shall be stored at
    \texttt{storage/models/unet\_model.h5}.
  \item The database shall be backed up to a volume snapshot before each snapshot
    submission.
\end{itemize}


\section{Legal and Ethical Considerations}
\label{sec:legal}


\begin{itemize}
  \item \textbf{ADA Compliance:} The system exists to support identification of
    sidewalk hazards that affect pedestrians with disabilities. Displacement
    measurements and flagging thresholds shall align with ADA accessibility standards.
  \item \textbf{Pedestrian Safety:} Rover operation in public spaces requires
    operator supervision at all times. The system shall not operate in autonomous
    mode without a qualified operator monitoring the live feed.
  \item \textbf{Image and Data Privacy:} Field scans may inadvertently capture images
    of pedestrians, vehicles, residential addresses, or license plates. All captured
    data shall be treated as sensitive, stored in access-controlled volumes, and not
    shared publicly or used for any purpose other than sidewalk assessment.
  \item \textbf{Data Retention:} Scan data shall be retained only as long as needed
    for engineering review and shall be deleted according to the LA City data
    retention policy.
  \item \textbf{Access Control:} Only authorized city engineering staff and project
    team members shall have access to the storage volume and database.
  \item \textbf{Open-Source Compliance:} All third-party libraries used in this
    project (TensorFlow, Open3D, FastAPI, etc.) shall be used in compliance with
    their respective open-source licenses.
\end{itemize}


\section{Glossary}
\label{sec:glossary}


\begin{longtable}{|p{1.5in}|p{5in}|}
\hline
\textbf{Term} & \textbf{Definition} \\
\hline
ADA & Americans with Disabilities Act. Federal law requiring accessible public
  infrastructure. \\
\hline
BOE & Bureau of Engineering, LA City department responsible for sidewalk
  assessment and remediation. \\
\hline
DEM & Digital Elevation Model. A raster image where pixel values represent
  ground elevation at that location. \\
\hline
GSD & Ground Sample Distance. The real-world distance represented by one pixel
  in the orthoimage. \\
\hline
LAS & A standardized binary file format for storing 3D point-cloud data from
  LiDAR sensors. \\
\hline
LiDAR & Light Detection and Ranging. A remote sensing technology that uses laser
  pulses to measure distance and generate 3D point clouds. \\
\hline
Orthoimage & A geometrically corrected aerial or top-down 2D image generated from a
  3D point cloud, with uniform scale. \\
\hline
PLY & Polygon File Format. An alternative point-cloud file format that stores
  geometry as a list of vertices and optionally faces. \\
\hline
ROS & Robot Operating System. A middleware framework providing libraries and tools
  for writing robot software. \\
\hline
U-Net & A convolutional neural network architecture designed for image segmentation
  tasks, used here to detect cracks and joints in sidewalk orthoimages. \\
\hline
Vertical Displacement & The signed elevation difference (in millimeters) between adjacent
  sidewalk panel segments at a joint or crack. Values above the threshold indicate
  a trip hazard. \\
\hline
Waypoint & A GPS coordinate (latitude and longitude) that the rover is instructed to
  navigate to during semi-autonomous operation. \\
\hline
\end{longtable}


\section{Future Work (Snapshot 4 Reflection)}


The following features were identified during the simulation but were not implemented and are
recommended for future development:

\begin{itemize}
  \item \textbf{Horizontal displacement calculation:} Research indicates that horizontal
    misalignment of sidewalk panels is an additional ADA concern. A paired-scan alignment
    method using ICP (Iterative Closest Point) is the recommended approach.
  \item \textbf{On-rover real-time inference:} Running the U-Net model directly on a
    Raspberry Pi 5 with a Neural Processing Unit (NPU) accelerator would eliminate the
    upload step and allow immediate field flagging.
  \item \textbf{Cloud storage and processing:} Migrating from local Docker volumes to
    Azure Blob Storage and Azure ML would improve scalability for city-wide deployments.
  \item \textbf{ROS2 migration:} The navigation module is designed around ROS Noetic
    (end-of-life 2025). A future iteration should migrate to ROS2 Humble.
  \item \textbf{YOLO-based rover centering:} A YOLO object detection model could center
    the rover on the sidewalk centerline automatically, reducing operator intervention
    during data collection.
  \item \textbf{Field validation at additional sites:} Testing at a wider variety of
    sidewalk materials, slopes, and urban density levels will improve model generalization.
\end{itemize}

\end{document}
